\scanpage{15}

Il ne semble pas qu'on puisse obtenir des normalisations locales sans \deleted{\unsure{type}} restriction de car. qui ferait intervenir moins de paramètres, [si on veut qu'elle soit subordonnée à la descente \deleted{\unsure{de}} par morphismes plats, ou \unsure{à|m.} simples]. Si on admet des ext. de base finies et plates, on peut pousser la normalisation plus loin en exigeant que $\delta(b,c,d,e,f)=1$ i.e. \inserted{(intrinsèquement)} se donner des trivialisations de $\unsure{\mathcal L|L}$ \deleted{compatibles} définies par des sections $\unsure{\lambda|D}$ dont la puissance $12$.ième (\unsure{ou même} $24$.ième) est \inserted{$\delta^{-1}$, où $\delta$ est} le discriminant. On va \deleted{\unsure{en ce}} donner des \deleted{\unsure{les}} normalisations plus poussées en car $\neq2$, puis en car. $\neq3$, puis on verra.
\mspage{14}
